% Generated draft: requires mathematical and source review.
\documentclass[11pt]{article}
\usepackage[margin=25mm]{geometry}
\usepackage{fontspec}
\setmainfont{Latin Modern Roman}
\usepackage{amsmath,amssymb,mathtools,unicode-math}
\setmathfont{Latin Modern Math}
\usepackage{enumitem,needspace,xurl,hyperref}
\setlength{\parindent}{0pt}
\setlength{\parskip}{5pt}
\setlength{\emergencystretch}{3em}
\newcommand{\N}{\mathbb{N}}
\newcommand{\Z}{\mathbb{Z}}
\newcommand{\Q}{\mathbb{Q}}
\newcommand{\R}{\mathbb{R}}
\newcommand{\C}{\mathbb{C}}
\newcommand{\F}{\mathbb{F}}
\newcommand{\A}{\mathbb{A}}
\newcommand{\PP}{\mathbb{P}}
\newcommand{\E}{\mathbb{E}}
\newcommand{\eps}{\varepsilon}
\newcommand{\dd}{\,\mathrm{d}}
\begin{document}
\section*{UnsolvedMath ID: 30006202}
\textbf{Source problem code:} catalogue entry
\subsection{Definitions and mathematical statement}
For a field \(k\), a polynomial \(k\)-algebra \(A\), and a group \(G\) acting on \(A\) by \(k\)-algebra automorphisms, the invariant-polynomial ring is
\[
A^G=\{f\in A: g\cdot f=f\text{ for every }g\in G\}.
\]
These data provide general notation only: the supplied record does not specify the field, polynomial algebra, class of groups, or permitted actions. A vector-space basis \(\mathcal B\) of \(A^G\) means a family such that every invariant polynomial has a unique expression as a finite \(k\)-linear combination of members of \(\mathcal B\). The record does not specify whether the intended bases are taken in the whole ring, in individual homogeneous components, or in another specified subspace.

\textbf{Research question, with unresolved qualifications.} For invariant-polynomial rings under various group actions, characterize approximately orthogonal bases, and determine when such bases can simultaneously satisfy the additivity property referred to in the source, analogous to that of tensor cumulants.

Orthogonality requires a specified inner product or other pairing; approximate orthogonality additionally requires a quantitative criterion and, if applicable, a limiting parameter and normalization. None is supplied. The original excerpt refers to an ``additivity property above,'' but that preceding property is absent from the record. Consequently, the operation, admissible inputs, independence assumptions if any, and identity defining additivity cannot be stated from the supplied evidence. In particular, additivity must not silently be interpreted as additivity under ordinary addition of polynomial arguments. The intended classes of actions and the quantifiers governing existence are also unspecified. Thus this is a statement of the supplied research direction with explicit missing data, rather than a fully quantified conjecture.
\subsection{Short English statement}
The problem asks what approximately orthogonal bases of invariant polynomials look like for general group actions, and when they can also have an additivity property analogous to tensor cumulants. The supplied excerpt omits the pairing, approximation criterion, and referenced additivity identity, so a precise existence statement requires further source context.
\subsection{Sources}
\begin{itemize}
\item Local UnsolvedMath registry, supplied complete source record: permanent ID \(30006202\), title \emph{Additive Orthogonal Bases for Invariant Polynomials}, exact source problem code \texttt{catalogue entry}. The record attributes the research direction to the Oberwolfach Reports workshop material from 2025; no individual proposer is supplied.
\item \emph{MFO-RIMS Tandem Workshop: Optimization, Theoretical Computer Science and Algebraic Geometry: Convexity and Beyond} (2025), Oberwolfach Reports, pp.\ 405--450, DOI \texttt{10.4171/owr/2025/9}. Supplied source and literature link: \url{https://doi.org/10.4171/owr/2025/9}. The supplied original excerpt carries the page marker 421 and refers to preceding material that is not included.
\item \emph{Tensor cumulants}. Supplied literature link: \url{https://arxiv.org/abs/2404.18735}. No author list or further bibliographic details are supplied.
\item Literature assessment checked at \texttt{2026-08-22}: the supplied curation labels the problem \texttt{partially\_solved} and states that tensor cumulants give a near-orthogonal additive basis for orthogonally invariant tensor polynomials, while the requested construction for general group actions remains unresolved. This is supplied dated curation derived from the OWR material, not an independently verified mathematical or current literature assessment. The supplied background requests independent verification of time-sensitive status and cited claims.
\end{itemize}
\end{document}
